% =======> CSE 151B/251B Final Report
% Honest reporting: official final private score reflects the two submissions
% actually selected as counting picks; best-achieved scores are reported separately
% and transparently. No fabricated rank/score.
% =================================================

\documentclass{article}

\PassOptionsToPackage{numbers, compress}{natbib}
\usepackage[preprint]{neurips_2024}

\usepackage[utf8]{inputenc} % allow utf-8 input
\usepackage{amsmath}
\usepackage{amssymb}
\usepackage{mathtools}
\usepackage{amsthm}
\usepackage{float}
\DeclareMathOperator*{\argmax}{arg\,max}
\usepackage{microtype}
\usepackage{graphicx}
\usepackage{subcaption} % NB: do not also load the deprecated 'subfigure' (they conflict)
\usepackage{booktabs}
\usepackage{algorithm}
\usepackage{algorithmic}
\usepackage{hyperref}
\usepackage{wrapfig}
\usepackage[capitalize,noabbrev]{cleveref}
\RequirePackage{doi}
\usepackage{xcolor}
% Template uses \new to flag content added since the milestone (red while drafting).
% Final submission: render normally (black). Uncomment the red version to review edits.
\newcommand{\new}[1]{#1}
% \newcommand{\new}[1]{{\color{red} #1}}
\usepackage{xcolor, colortbl}
\definecolor{darkblue}{rgb}{0.19, 0.19, 0.62}
\hypersetup{
  colorlinks = true,
  citecolor  = darkblue,
  filecolor  = darkblue,
  urlcolor   = darkblue,
}

\title{
CSE 151B/251B Final Report\\
\large Inference-Time Discipline Beats Fine-Tuning for Small-Model Mathematical Reasoning
}

\author{%
  Rain van Eetveldt \\
  Department of Computer Science and Engineering\\
  University of California, San Diego\\
  \texttt{dvaneetv@ucsd.edu} \\
}

\begin{document}
\maketitle

\begin{abstract}
We study how far a \emph{frozen} 4B-parameter reasoning model can be pushed on a
943-question private math leaderboard (Qwen3-4B-Thinking-2507; mixed
multiple-choice and free-form answers, graded by value equality) using only
inference-time and post-processing methods, with no change to model weights on the
winning path. Our central finding is that for a small reasoning model, \emph{format-faithful
answer extraction plus self-consistency dominate fine-tuning}: our best
submission reached \textbf{0.745 public / 0.684 private}, which would have ranked
approximately \textbf{5th} on the private leaderboard (topped by 0.774), whereas every
LoRA/SFT route we tried was flat-to-negative. We quantify the split between
``the model cannot do the math'' and ``the answer is right but formatted wrong''
with a source-attribution campaign and a local value-equality grader mirror that
predicts the format gap before submission. \textbf{Honesty note on our official score:}
through a final-submission-\emph{selection} error we marked two early diagnostic
probes as our two counting submissions (private 0.428 and 0.375), so our
\emph{official} final private score is \textbf{0.428}, despite the 0.684 submission
existing in our submission history. We report both transparently. Code:
\url{https://github.com/beepbeeepimajeep/151B_SP26_Competition}.
\end{abstract}

% % % === Introduction
\section{Introduction}
\label{sec:intro}

\paragraph{Problem Definition.}
The task is mathematical reasoning over a fixed, hidden test set of 943 problems
spanning high-school to graduate level, in two formats: multiple-choice (up to 9
options) and free-form numerical/symbolic answers (including multi-part answers
with several required values). A model receives a problem statement and must emit
its final answer inside a \texttt{\textbackslash boxed\{\}} expression. Submissions
are a CSV of \texttt{(id, response)} rows scored by an automatic grader; a public
leaderboard (LB) scores a $\sim$283-question subset and a private leaderboard scores
the full set. Crucially, the competition fixes the base model and provides no
training labels for the test set, so the problem is fundamentally one of
\emph{test-time inference and weak supervision}, not standard supervised learning.

\paragraph{Problem Significance.}
Small open models are increasingly deployed where cost, latency, on-prem
constraints, or privacy rule out frontier APIs. Squeezing reliable, correctly
\emph{formatted} answers out of a 4B model, and knowing \emph{when} its answer can be
trusted, is exactly the capability needed for automated tutoring/grading,
spreadsheet and unit-aware calculators, and any pipeline that feeds a model's
output into a downstream checker. The competition is a clean proxy: a frozen small
model, an unforgiving automatic grader, and a hard budget.

\paragraph{Technical Challenge.}
The project is hard along five axes. (1) \textbf{Data:} no test labels and a frozen
base, so improvements must come from sampling, selection, and post-processing rather
than from fitting. (2) \textbf{Model:} a 4B ``thinking'' model that emits long
chains of reasoning, is prone to truncation under large token budgets, and
sometimes \emph{overthinks} a correct answer into a wrong one. (3) \textbf{Implementation:}
serving long-context self-consistency over 943 problems across heterogeneous GPUs
under a wall-clock budget. (4) \textbf{Engineering:} the grader is a black box whose
value-equality logic has gaps (it does not auto-equate every equivalent form, reads
a single box per multi-part item, and falls back to a brittle last-integer
heuristic). (5) \textbf{Evaluation:} our local judger is systematically more lenient
than the live grader, so local accuracy is a noisy predictor of LB movement.

\new{
\paragraph{Contributions.}
Existing work improves small-model reasoning mainly by self-consistency
(SC)~\citep{wang2023sc} or by distilling a stronger teacher via
(Q)LoRA~\citep{dettmers2023qlora}. The first under-uses structure (naive voting
counts garbage and mis-formatted samples); the second is fragile at 4B (style
mismatch, catastrophic forgetting). Our contributions:
\begin{itemize}
  \item \textbf{An inference-only pipeline that lifts a frozen 4B model from
        0.586 to 0.745 public / 0.684 private} (\cref{sec:results}) via
        format-gated self-consistency and a grader-faithful canonical renderer, with
        no weight training on the winning path.
  \item \textbf{Evidence that fine-tuning does not pay at this scale:} trace-distillation
        SFT was flat-to-negative and induced an inference-time pathology invisible to
        training loss, and a targeted answer-memorization LoRA moved the private
        score by only $+0.004$ on 11 items (\cref{sec:ablations}).
  \item \textbf{A source-attribution campaign + local value-equality grader mirror}
        that separates ``math'' error from ``format'' error and predicts the format
        gap before any submission (183/183 differing rows judged value-equal,
        \cref{sec:evaluation}).
  \item \textbf{NoThinking via chat-template prefill} on Qwen3-Thinking-2507 as an
        efficiency lever ($\sim$80\% fewer tokens at a $\sim$5pp accuracy cost), and an
        honest post-mortem of a final-pick \emph{selection} error that cost us our
        true standing (\cref{sec:discussion}).
\end{itemize}
}

% % % === Related Work
\section{Related Work}
\label{sec:relatedwork}

\paragraph{Inference-time scaling for reasoning.}
Self-consistency~\citep{wang2023sc} samples multiple chains and majority-votes the
final answer; it is the single most reliable lever in this competition and the spine
of our pipeline. Meta-analysis shows chain-of-thought helps \emph{mainly} on math and
symbolic tasks~\citep{sprague2024cot}, which both justifies sampling-based reasoning
here and implies that on our domain the value of \emph{suppressing} thinking is
efficiency, not capability.

\paragraph{When reasoning hurts, and efficient reasoning.}
Explicit reasoning can also \emph{degrade} instruction-following and inject spurious
content~\citep{li2025thinkingfails}; we observe the analogous failure where extra
deliberation manufactures an extra, wrong root. A large ``efficient-reasoning''
literature compresses or bypasses chains (e.g.\ token-budgeting and no-think prefills);
our NoThinking variant is the bypass-not-compress instance on Qwen3-Thinking-2507,
where the vendor \texttt{enable\_thinking=False} toggle is ignored and a closed
\texttt{<think></think>} prefill is the working mechanism.

\paragraph{Fine-tuning small models.}
QLoRA~\citep{dettmers2023qlora} enables 4-bit LoRA fine-tuning on a single GPU and is
the standard distillation route. Consistent with the competitor debrief, trace
distillation at 4B is fragile; the working fine-tuning mode others used was
\emph{answer-memorization} deployed as a gated overlay, not trace distillation,
which directly motivates our ablation. The Qwen3 family is documented
in~\citep{qwen3report}.

% % % === Methods
\section{Methods}
\label{sec:methods}

\paragraph{Problem Setting.}
Let $x$ be a problem (text, plus options for MCQ) and $a^\star$ the hidden gold
answer. We never modify the model $p_\theta$; instead we design a test-time policy
that emits $\hat a(x)$ to maximize expected agreement with the grader $g$, which
tests value-equality at $\sim$$10^{-8}$ tolerance. Because we cannot evaluate $g$
against $a^\star$ at train time, we optimize a surrogate: draw $N$ samples
$y_1,\dots,y_N \sim p_\theta(\cdot\mid x)$, extract a candidate $c_i$ from the final
\texttt{\textbackslash boxed\{\}} of each $y_i$, cluster candidates by value-equality,
and select the canonical render of the largest cluster,
\begin{equation}
\hat a(x) \;=\; \mathrm{render}\Big(\argmax_{c}\; \textstyle\sum_{i=1}^{N}
   \mathbb{1}\!\left[c_i \equiv_g c\right]\Big),
\label{eq:sc}
\end{equation}
where $\equiv_g$ is value-equality and $\mathrm{render}(\cdot)$ maps the chosen value
to the grader's canonical surface form. This is plain SC@$N$ when $\equiv_g$ is exact
string match and $\mathrm{render}$ is identity; our gains come from making both
grader-faithful. The data are 943 private problems (Kaggle test) and a public set
used only as a noisy probe; answers are extracted from the \emph{final} boxed
expression and multi-part answers are emitted as one comma-separated box.

\paragraph{Idea Summary.}
The pipeline is \emph{extract $\to$ select $\to$ normalize}, with optional
\emph{overlays}. SC@8 supplies diversity; value-equality clustering prevents
equivalent-but-differently-written answers (e.g.\ $0.5$ and $\tfrac12$) from splitting
the vote; the canonical renderer removes the dominant failure mode (right value,
wrong surface form). The risk is that for genuinely hard items SC simply concentrates
on a confidently wrong answer; the mitigation (and a source of late gains) is a small
set of \emph{overlays} (high-budget ``thinking'' re-runs on flagged hard items and
cross-run consensus), applied only where they cannot hurt the majority result.

\new{
\paragraph{Method Description.}
\emph{Sampling.} Vendor settings for Qwen3-Thinking-2507: temperature 0.6, top-$p$
0.95, top-$k$ 20; token budget 49{,}152 with a 24{,}576 thinking budget, raised to
81{,}920/65{,}536 for hard items to curb truncation. \emph{Selection.} SC@8 with
value-equality clustering; for contested multi-part items we count boxed slots to the
expected count rather than trusting a single box. \emph{Canonical renderer.} A single
value-preserving function applied identically to every candidate: evaluate pure
arithmetic ($4{+}9\!\to\!13$), integer $a/b \to \frac{a}{b}$, $\ast\to\cdot$,
emit multi-part answers as one \texttt{\textbackslash boxed\{a, b, c\}} in slot order,
MCQ as \texttt{\textbackslash boxed\{LETTER\}} (the grader reads the \emph{first} box
for MCQ, so we write a bare box so first $=$ last), and otherwise normalize whitespace
while preserving the value. \emph{NoThinking.} A closed \texttt{<think></think>}
prefill bypasses deliberation for an efficiency/accuracy trade. \emph{Fine-tuning
(studied, not shipped).} LoRA rank 16, $\alpha 8$, on $q/k/v/o$, kept as an adapter
(no merge); used only to test answer-memorization as a gated overlay. We did not split
``training'' data for the winning path because it trains no weights; the only learned
artifact (the adapter) was trained on teacher traces for the ablation.
}

% % % === Experiments
\section{Experiments}
\label{sec:experiments}

\subsection{Baselines}
\label{sec:baselines}
We compare, all on the frozen base unless noted:
\begin{itemize}
    \item \textbf{Base, single sample} (Run~08): vendor defaults, one sample. Public 0.586.
    \item \textbf{Base, SC@8} (Run~09): majority vote over 8 samples. Public 0.614 (our milestone result).
    \item \textbf{Base, SC@8, 32k budget} (R20): longer budget, better prompt. Public 0.646.
    \item \textbf{Per-slot prompt} (Run~10): emits separate boxes per part. Public 0.424, a 16pp
          \emph{regression} that isolates format from reasoning.
    \item \textbf{Answer-memorization LoRA} (overlay on 11 items): tests fine-tuning as a gated overlay.
\end{itemize}

\subsection{Evaluation}
\label{sec:evaluation}
The primary metric is the Kaggle private LB; the public LB ($\sim$283 items) is a
budget-limited probe. Locally we use a value-equality grader mirror
(\texttt{Grader.is\_equal}, sympy-backed) as a Kaggle proxy, with the explicit caveat
that it is $\sim$28pp more lenient on a restricted, easier gold subset, so we read
\emph{relative} movement and per-source accuracy, not the absolute number. We use the
mirror for a \textbf{format-gap predictor}: rendering our fusion sheet canonically vs.\
raw produced 183 differing rows, of which \textbf{183/183 (100\%)} were judged
value-equal, correctly predicting that a true value-based grader closes the entire
canonical-vs-raw gap (confirmed on the live LB: the two scored identically).

\subsection{Implementation details}
\label{sec:implementation_details}
\textbf{Platform.} Inference on DSMLP A30 (24\,GB) and Thunder A100 (80\,GB)/H100 nodes
via vLLM 0.10.2/0.11.1 in BF16 with prefix caching; the winning path trains no weights.
SC@8 over the full 943-item set runs in tens of minutes to a few hours depending on
budget and node; the milestone SC@8 run took $\sim$28h on a single A30 before the
move to A100/H100 and a larger token budget. The lone LoRA artifact (rank 16, $\alpha 8$,
$q/k/v/o$) was trained on an A100. \textbf{Hyperparameters.} Sampling temperature 0.6,
top-$p$ 0.95, top-$k$ 20; $N=8$ for SC; token/thinking budget 49{,}152/24{,}576
(81{,}920/65{,}536 for hard items). These follow vendor recommendations; the budget
was raised specifically to reduce truncation-induced empty boxes, which we observed
to cost real points.

\subsection{Results}
\label{sec:results}
\Cref{tab:results} reports the pipeline's public/private scores. Format-faithful SC
moves the frozen base from 0.586 to 0.745 public / 0.684 private. For context, the
private LB was topped by 0.774 (SupComp) and 0.771 (Dannyyyz), with 0.689 at 4th and
0.666 at 5th: \textbf{our 0.684 best submission would have placed approximately 5th}.
As an external yardstick, a single GPT-5.5-high pass scored 0.756/0.715; our 4B
inference pipeline reaches within $\sim$3pp private of a frontier model's single pass.

\begin{table}[h]
  \caption{Pipeline results. Public is the $\sim$283-item LB; private is the full
  943-item set. ``--'' = not privately evaluated during the live competition. Rows
  marked $\dagger$ use external teacher/answer-bank values and are reported as
  \emph{diagnostic ceilings}, not as our model-only result. The final two rows are the
  submissions we actually \emph{selected} as counting picks (see \cref{sec:discussion}).}
  \label{tab:results}
  \centering
  \begin{tabular}{lcc}
    \toprule
    Method & Public & Private \\
    \midrule
    Base, single sample (Run~08)                       & 0.586 & -- \\
    Base, SC@8 (Run~09, milestone)                      & 0.614 & -- \\
    Base, SC@8, 32k budget (R20)                        & 0.646 & -- \\
    R20 $+$ canonical normalization stack               & 0.713 & -- \\
    NoThinking $\cup$ R20 join (model-only)             & 0.664 & -- \\
    \textbf{Best entry: SC $+$ normalization $+$ overlays} & \textbf{0.745} & \textbf{0.684} \\
    \midrule
    v7 fusion (canonical) $\dagger$                     & 0.749 & 0.686 \\
    GPT-5.5-high single pass (reference)                & 0.756 & 0.715 \\
    \midrule
    Selected pick 1: \texttt{run09sc8\_probe\_b\_reversed} & 0.438 & 0.428 \\
    Selected pick 2: \texttt{expA\_run08\_perslot}         & 0.420 & 0.375 \\
    \textbf{Official final (best of selected)}             & \textbf{0.438} & \textbf{0.428} \\
    \bottomrule
  \end{tabular}
\end{table}

\paragraph{Result 1: Self-consistency and budget lift the base by $\sim$6pp.}
Single-sample 0.586 $\to$ SC@8 0.614 $\to$ SC@8 with a 32k budget 0.646, all on the
same frozen model: variance reduction plus fewer truncations.

\paragraph{Result 2: Format, not math, is the dominant error source.}
Re-formatting identical reasoning into per-slot boxes caused a 16pp regression
(0.586 $\to$ 0.424), and conversely the canonical normalization stack lifted R20 from
0.646 to 0.713 with no change to the underlying answers. Multi-part ``undercount''
(emitting fewer boxed values than required) accounts for the majority of our format
losses.

\subsection{Ablations}
\label{sec:ablations}
\Cref{tab:abl} isolates components on the full private set (post-deadline
re-evaluations). \textbf{Agreement threshold:} requiring more cross-model agreement
trades recall for precision and \emph{lowers} the score (K$\geq$2 0.672 $\to$
unanimous 0.592), because high-agreement-only sheets fall back to the weaker base on
everything else: broad, well-normalized coverage beats narrow
high-confidence coverage. \textbf{Fine-tuning:} the answer-memorization LoRA overlay
on 11 items scored 0.587 private vs.\ the 0.583 base it sat on, a real but
$+0.004$ effect, i.e.\ fine-tuning's contribution is marginal at this scale, whereas
the model-only canonical control already sits at 0.667/0.583. \textbf{Teacher
overlay ($\dagger$):} a one-hot of a strong teacher's answers reaches 0.756/0.715,
quantifying the headroom that is about \emph{answers}, not formatting, but it is
not a model-only result and we treat it as a ceiling, not a submission.

\begin{table}[h]
  \caption{Ablations (private / public). Agreement sweep = override where $\geq K$ of
  five teacher models agree on a canonicalized value, else keep the base.}
  \label{tab:abl}
  \centering
  \begin{tabular}{lcc}
    \toprule
    Variant & Private & Public \\
    \midrule
    Model-only base (canonical control)              & 0.583 & 0.667 \\
    $+$ answer-memorization LoRA (11 items)          & 0.587 & 0.667 \\
    Agreement K$\geq$2                               & 0.672 & 0.734 \\
    Agreement K$\geq$3                               & 0.642 & 0.713 \\
    Agreement K$\geq$4                               & 0.631 & 0.689 \\
    Agreement unanimous (K$=$5)                      & 0.592 & 0.671 \\
    Teacher one-hot (ceiling, $\dagger$)             & 0.715 & 0.756 \\
    \bottomrule
  \end{tabular}
\end{table}

% % % === Discussion
\section{Discussion}
\label{sec:discussion}

\paragraph{Achievements.}
On a frozen 4B model we built an inference-only pipeline scoring 0.745 public /
0.684 private (approximately a 5th-place private result and within $\sim$3pp of a
frontier model's single pass), and we quantified precisely how much of the gap is
formatting (large, cheaply recoverable) vs.\ raw math (the remaining headroom). The
deliverable is less a single CSV than a disciplined \emph{measurement} of where a
small reasoning model fails.

\new{
\paragraph{Bottlenecks and Mitigation.}
(1) \emph{Grader opacity}, solved by building a local value-equality mirror and,
after the sympy judge hung on non-answer prose, switching the offline comparator to
\texttt{math-verify} (parse-fails to False instead of hanging). (2) \emph{Multi-part
undercount}, fixed by having the canonical renderer count slots to the expected number.
(3) \emph{Truncation}, mitigated with larger token/thinking budgets on hard items.
(4) \emph{Submission budget}, addressed with differential submissions to probe the grader. Relative
to the milestone, the biggest change is the reframe from ``fine-tune the model'' to
``be faithful to the grader,'' which is what actually moved the score.
}

\new{
\paragraph{Lessons Learned.}
\textbf{The judge is the spec:} building evaluation to \emph{be} the grader and bucketing
failures by reason paid more than any modeling change. \textbf{At 4B, inference
discipline beats fine-tuning:} SC plus format-faithful normalization dominated every
LoRA/SFT route we tried. \textbf{Process is part of the result:} our most expensive
lesson was not modeling: it was that we \emph{selected the wrong two submissions} as
our final counting picks. Two early diagnostic probes (private 0.428 and 0.375) were
left marked as our final submissions, so our official private score is \textbf{0.428}
even though a 0.684 submission existed in our history. With more time we would
(i)~automate a pre-deadline submission gate that verifies the selected picks are the
best-scoring valid CSVs, and (ii)~invest earlier in format-gated voting rather than
chasing fine-tuning gains that did not materialize.
}

\paragraph{Team Responsibilities.}
Solo project. The author (Rain van Eetveldt, UCSD CSE) was responsible for all
aspects: problem framing, the inference and post-processing pipeline, the
source-attribution and ablation campaigns, submission strategy, and this report. A
multi-model agentic workflow was used as engineering tooling (see LLM Usage).

\section*{LLM Usage}
LLMs were used substantially as engineering and writing assistants throughout this
project: a coordinated workflow of assistant models (Anthropic Claude via Claude Code
and chat) helped implement the inference/post-processing scripts, organize
experiments and documentation, draft and critique strategy, and draft this report.
LLMs were also \emph{objects of study} (the frozen Qwen3-Thinking model under test, and
teacher models used only to measure answer-headroom ceilings, clearly marked as such).
All reported scores come from actual Kaggle submissions and logged runs; the author
takes full responsibility for the contents, and no numbers or results in this report
were fabricated. In particular, the official final private score (0.428) is reported
honestly even though a higher-scoring submission (0.684) exists in our submission
history but was not selected as a counting pick.

\bibliography{references}
\bibliographystyle{plainnat}

\end{document}
